\begin{abstract}
Skills supplied in a language model's context can improve task performance,
but how the model represents and uses their content remains unclear. Can one
hidden-state vector carry a skill's effect, or are states at multiple token
positions needed? We compare a correct skill with a wrong skill of the same
token length, then inject the difference between their hidden states. In
synthetic tasks, the final prompt position carries a large change shared by
both skills, yet the content-specific difference accounts for most of the
behavioural gain. Recovery at this position depends on answer format: it falls
from $1.000$ for a one-token answer to $0.235$ for a short number and $0.003$
for step-by-step reasoning. Injecting states across the skill's own token span
instead recovers $0.89$ of its gain on a screened clinical-calculator subset
of $135$ items. Recovery depends on both skill content and token order.
Keeping only the leading $32$ directions of the centred state difference
recovers $0.12$, compared with $0.77$ for $128$ directions. The rank needed to
reach half of the measured recovery range does not scale proportionally with
model width within Qwen; Mistral's is higher, though with overlapping intervals. Across two task
domains and two model families, paired interventions further show that span
states stop transferring most of the skill's effect before the model stops
depending on attention to them. These findings identify a transferable
representation across skill tokens and clarify the limits of probing it at a
single prompt position.
\end{abstract}
